%=============================================================================
% OSR POSITIONING — UIDT v3.9
% Explicit placement within the current Ontic Structural Realism landscape.
% TKT-2026-05-12-manuscript-elite-revision Step 3/4
%
% Addresses review finding: UIDT needs explicit comparison with
% dynamical/structural-actualist OSR variants (Ladyman/Ross, French,
% Esfeld/Deckert, McKenzie) to achieve maximum community acceptance.
%=============================================================================
\subsection{\textsc{uidt} as Dynamically Grounded Ontic Structural Realism}
\label{sec:osr_positioning}

The philosophical position of \textsc{uidt} can be situated precisely within
the contemporary landscape of ontic structural realism (OSR). The central
thesis of OSR---that the fundamental furniture of the world consists of
\emph{relations and structures} rather than intrinsically characterised
objects~\cite{Ladyman2007,French2014,McKenzie2017}---is fully consonant
with the \textsc{uidt} identification of the vacuum information-density
field $S(x)$ as the carrier of physical structure, with particles and
fields emerging as stable relational patterns therein.

However, not all OSR variants are equivalent, and \textsc{uidt} aligns
most closely with \emph{dynamically grounded} or \emph{minimalist} variants
rather than with purely formal-structural readings. Specifically:

\begin{enumerate}[leftmargin=1.8em,itemsep=4pt]
  \item \textbf{Against eliminativist OSR (French~\cite{French2014}).}
        \textsc{uidt} does not eliminate objects entirely; it
        \emph{grounds} them in the dynamics of $S(x)$. Individual
        particle-like excitations are emergent stable configurations of
        the vacuum lattice, not mere notational conveniences.
        This places \textsc{uidt} closer to \emph{moderate} or
        \emph{entity-in-context} OSR.
  \item \textbf{Affinity with minimalist ontology (Esfeld and
        Deckert~\cite{Esfeld2017}).}
        The \textsc{uidt} vacuum lattice shares the spirit of Esfeld and
        Deckert's minimalist ontology: a sparse set of fundamental entities
        (information-density nodes) whose intrinsic properties are minimal,
        with all physical content encoded in their relational configuration.
  \item \textbf{Affinity with structural actualism
        (McKenzie~\cite{McKenzie2017}).}
        \textsc{uidt} treats the Yang--Mills symmetry group $\mathrm{SU}(3)$
        not as a descriptor of a pre-existing ontology but as the
        constitutive constraint on what configurations of $S(x)$ can be
        physically realised. This resonates with structural actualism's
        insistence that modal structure is grounded in concrete physical
        dynamics.
  \item \textbf{Departure from purely informational idealism.}
        Unlike some information-theoretic programmes that treat the wavefunction
        or information as epistemically or subjectively grounded, \textsc{uidt}
        assigns $S(x)$ an ontic status: it is a real field with a Lagrangian,
        coupling constants, and measurable consequences. This distinguishes
        \textsc{uidt} from QBist accounts~\cite{Fuchs2014} at the ontological
        level, even while sharing their epistemic stratification methodology.
\end{enumerate}

\begin{tcolorbox}[pillarbox,
  title={\small Positioning Summary: \textsc{uidt} and OSR Variants},
]
{\small
\begin{tabular}{@{}p{4.2cm}p{3.0cm}p{5.2cm}@{}}
  \toprule
  \textbf{OSR Variant} & \textbf{Alignment} & \textbf{Key difference} \\
  \midrule
  Eliminativist (French) & Partial & UIDT retains emergent individuation \\
  Minimalist (Esfeld/Deckert) & Strong & Sparse nodes + relational dynamics \\
  Structural actualism (McKenzie) & Strong & SU(3) as constitutive constraint \\
  QBism (Fuchs) & Epistemic only & $S(x)$ is ontic, not agent-relative \\
  Mathematical universe (Tegmark) & Partial & UIDT is physically motivated, not pan-mathematical \\
  \bottomrule
\end{tabular}
}
\end{tcolorbox}

\noindent
This positioning is not merely taxonomic. It has a direct methodological
consequence: \textsc{uidt} is falsifiable at the \emph{structural} level
(if the $\mathrm{SU}(3)$ coupling does not produce the predicted gap, the
entire information-density ontology loses its primary physical anchor),
whereas purely formal-structural accounts resist falsification at the
ontological level. The explicit Kill-Switch criteria (F1--F7) in
Section~\ref{sec:killswitches} reflect this commitment to an empirically
responsible structural realism.
